\begin{bild}{Integrierte Ergänzung v1.6: dieselbe Hülle als Matrixordnung}
Neben der $D_5$--$A_3$-Verklebung existiert eine ausdrücklich markierungstreue Darstellung als maximale $\Z[i]$-Ordnung $M\subset M_2(\mathbb Q(i))$. Mit $Q(A)=\operatorname{Re}\Tr(AA^\dagger)$ ist ihr achtdimensionales ganzzahliges Gitter positiv, gerade und unimodular; es hat genau 240 Vektoren mit $Q=2$. Der vorhandene Compiler wird durch eine ganzzahlige Isometrie samt komplexer Struktur und Familienwirkung getroffen. Das ist keine Gleichsetzung von 240 Koordinaten mit 240 Wurzeln, sondern eine geprüfte Abbildung. Die Wahl der maximalen statt einer kleineren Ordnung bleibt eine zusätzliche Frage. Der vollständige Beweis und seine Herkunft stehen in Kapitel~\ref{ch:v16-matrix}.
\end{bild}
